\subsection{Reconstruction from compatible restrictions}

Finite windows determine a single graded object only after their
restrictions have been identified.  A \emph{compatible tower} consists
of bialgebra objects $C^{(N)}$ in
$\mathsf{Gr}_{\leq N}(Q_+;\mathcal V)$ and supplied bialgebra
isomorphisms
\[
 \theta_N:\tau_{\leq N}C^{(N+1)}\xrightarrow{\sim}C^{(N)}
 \qquad(N\geq0).
\]
Their composites define
$\theta_{M,N}:\tau_{\leq M}C^{(N)}\xrightarrow{\sim}C^{(M)}$
for $M\leq N$, with $\theta_{N,N}=1$.
They satisfy
\[
 \theta_{L,N}=\theta_{L,M}\,\tau_{\leq L}(\theta_{M,N})
 \qquad(L\leq M\leq N).
\]
A morphism of towers $f:C^{(\bullet)}\to D^{(\bullet)}$ is a family
of bialgebra maps $f_N:C^{(N)}\to D^{(N)}$ satisfying
\begin{equation}
 f_N\theta_N^C=\theta_N^D\,\tau_{\leq N}(f_{N+1}).
 \label{tower-morphism}
\end{equation}

A \emph{complete graded bialgebra} here has a specified product
decomposition $\widehat C=\prod_{\beta\in Q_+}C_\beta$, the height
topology and tensor product~\eqref{completed-tensor}.
Its multiplication, unit, coproduct and counit are continuous and
preserve total charge.  Morphisms are continuous charge-preserving
bialgebra maps.  These definitions use the chosen ordinary, super or
cochain base category, with its stated symmetry.

\begin{theorem}[Reconstruction from compatible windows]\label{tower-theorem}
The following categories are equivalent:
\begin{enumerate}
\item nonnegatively $Q_+$-graded bialgebras in $\mathcal V$;
\item compatible towers with the supplied restriction isomorphisms;
\item complete graded bialgebras with the specified product decomposition,
height topology and tensor product~\eqref{completed-tensor}.
\end{enumerate}
In the cochain case, products and inverse limits are taken in each
cohomological degree.  No finite-dimensionality hypothesis is required.
\end{theorem}

\begin{proof}
A graded bialgebra gives a tower by Proposition~\ref{monoidal-restriction}.
Conversely, for $h=\htop(\beta)$ put $C_\beta=C^{(h)}_\beta$.
For $N\geq h$, write
\[
 u_{\beta,N}=(\theta_{h,N})_\beta:C^{(N)}_\beta
          \xrightarrow{\sim}C_\beta.
\]
For $\gamma+\delta=\beta$ and $N\geq\htop(\beta)$ define
\begin{align*}
 m_{\gamma,\delta}
 &=u_{\beta,N}m^{(N)}_{\gamma,\delta}
           (u_{\gamma,N}^{-1}\otimes u_{\delta,N}^{-1}),\\
 \Delta_{\gamma,\delta}
 &=(u_{\gamma,N}\otimes u_{\delta,N})
           \Delta^{(N)}_{\gamma,\delta}u_{\beta,N}^{-1}.
\end{align*}
The isomorphism $\theta_N$ preserves both operations, and
$u_{\beta,N+1}=u_{\beta,N}(\theta_N)_\beta$.
Substitution in either formula cancels the two adjacent $\theta_N$
factors.  The value at stage $N+1$ therefore equals the value at stage
$N$.  Induction proves independence of every later stage.

Transport the unit by $u_{0,N}\eta^{(N)}$, and the charge-zero counit by
$\varepsilon^{(N)}u_{0,N}^{-1}$.  The same cancellation proves their
independence.  The counit is zero in every positive charge.
In the cochain case transport the differential by
$u_{\beta,N}d^{(N)}u_{\beta,N}^{-1}$.  Each $\theta_N$ is a chain map,
so this too is stage independent, has degree $+1$ and squares to zero.
In the super case the transports preserve parity.

The maps above define multiplication and coproduct on
$C=\bigoplus_\beta C_\beta$.  For each $\beta$, only finitely many
pairs $\gamma+\delta=\beta$ occur, so the coproduct of an element has
finite support.  Each associativity, coassociativity or compatibility
identity has a fixed total charge.  Choose a stage containing that
charge and transport its identity to $C$.
This proves all identities in~\eqref{bialgebra-identities}, as well as
the algebra and coalgebra axioms.  It also proves both unit and counit
comparisons and all chain-map identities, including the derivation and
coderivation equations.

For a tower morphism, set
\[
 f_\beta=u^D_{\beta,N}(f_N)_\beta(u^C_{\beta,N})^{-1}
 \qquad(N\geq\htop(\beta)).
\]
Equation~\eqref{tower-morphism} proves independence of $N$ by the same
adjacent-stage cancellation.  Transport of the structure-map squares
shows that $\bigoplus_\beta f_\beta$ is a bialgebra map, and a chain map
when appropriate.  This construction preserves identities and
composition componentwise.  The maps $u_{\beta,N}$ assemble to natural
bialgebra isomorphisms
$C^{(N)}\xrightarrow{\sim}\tau_{\leq N}C$ commuting with the
$\theta_N$.  Starting from a graded object recovers its original
components and maps.  This proves the equivalence of the first two
categories, on objects and morphisms.

On $\widehat C=\prod_\beta C_\beta$ define
\begin{equation}
 (ab)_\beta=\sum_{\gamma+\delta=\beta}
 m_{\gamma,\delta}(a_\gamma,b_\delta),\qquad
 (\widehat\Delta a)_{\gamma,\delta}
   =\Delta_{\gamma,\delta}(a_{\gamma+\delta}).
 \label{reconstructed-operations}
\end{equation}
The multiplication extends linearly from the completed tensor product
by the same finite sum in each charge.  The unit and counit use charge
zero, and the differential acts componentwise.
Each output through total height $N$ depends only on input components
through height $N$.  Thus these linear structure maps are continuous.
For multiplication of elements, the subspaces of heights greater than
$N$ are two-sided ideals, which also proves joint continuity.
Every identity holds after projection to each total-height window,
hence in the separated completion.  The tensor associators and symmetry
extend for the same reason.

Conversely, restrict the structure maps of a complete graded bialgebra
to homogeneous components.  Total-charge preservation and the finite
number of decompositions of a charge give a graded bialgebra on their
direct sum.  That direct sum is dense in the product, so continuity
recovers exactly the original maps by~\eqref{reconstructed-operations}.
For morphisms, taking the product of the reconstructed component maps
is continuous and preserves every operation.  Restriction and completion
are inverse on these maps and preserve identities and composition.
This proves the third equivalence.
\end{proof}

The supplied isomorphisms are essential.  The existence of a bialgebra
in each height does not identify its restriction with the preceding one.
Nor does the theorem identify an arbitrary complete topological object
with a product of prescribed homogeneous pieces.

Finite-dimensional homogeneous vector spaces give finite-dimensional
windows because there are only finitely many charges below each bound.
Their completion is then linearly compact over the discrete field.
Neither property is needed above.  For complexes, degreewise finiteness
means $\dim C_\beta^j<\infty$ for each $j$.
It gives finite-dimensional windows in each cohomological degree, but
need not give finite total dimension or bounded complexes.
The completion satisfies $\widehat C^j=\prod_\beta C_\beta^j$.
It imposes no completion in cohomological degree.

\subsection{Rectangular coproduct targets}

A coproduct with separate bounds on its two outputs retains more terms
than a tensor window with one total bound.  The completed tensor also
records these separate bounds.

\begin{proposition}[Ordinary tensor targets]\label{rectangular-tensor}
There is a canonical identification
\[
 \widehat V\widehat\otimes_{\mathrm{ht}}\widehat W
 \cong\varprojlim_{r,s}
 \bigl(\tau_{\leq r}V\otimes\tau_{\leq s}W\bigr).
\]
The finite-output coproduct of a reconstructed bialgebra has the type
\begin{equation}
 \Delta_{r,s}:\tau_{\leq r+s}C
 \longrightarrow\tau_{\leq r}C\otimes\tau_{\leq s}C.
 \label{rectangular-coproduct}
\end{equation}
For $r\leq R$ and $s\leq S$, write $p_a^b$ for height projection
from bound $b$ to bound $a$.  Then
\begin{equation}
 (p_r^R\otimes p_s^S)\Delta_{R,S}
   =\Delta_{r,s}p_{r+s}^{R+S}.
 \label{rectangular-restriction}
\end{equation}
\end{proposition}

\begin{proof}
Triangular and rectangular height regions are cofinal, since
\[
 \{i+j\leq N\}\subseteq\{i,j\leq N\}
 \subseteq\{i+j\leq2N\},
\]
and any rectangle with bounds $r,s$ lies in the triangle of bound $r+s$.
Both inverse limits consequently specify one element of
$V_\gamma\otimes W_\delta$ for every pair of charges.
Each finite region has finitely many charge pairs.  Tensor products
distribute over these finite sums without any dimension hypothesis.
For complexes this argument applies in each cohomological degree,
using the ordinary tensor complex inside each charge pair.

The map $\Delta_{r,s}$ retains precisely those components
$\Delta_{\gamma,\delta}$ with
$\htop(\gamma)\leq r$ and $\htop(\delta)\leq s$.
Their input charge is $\gamma+\delta$, of height at most $r+s$.
Equation~\eqref{rectangular-restriction} retains the same components
in either order.  Equivalently, the rectangular projection of
$\widehat\Delta$ factors through input height $r+s$.
\end{proof}

The polynomial formula~\eqref{polynomial-cross-term} shows why this
input bound matters.  With $r=s=1$, the output retains $2x\otimes x$
and therefore needs the charge-two input $x^2$.

\subsection{Antipodes}

An antipode is a two-sided inverse of the identity under convolution.
The connected graded construction and its finite sum formula are given
in~\cite[Propositions 1.4.14 and 1.4.22, pp.~15, 17]{GrinbergReiner}.
The height grading makes that construction compatible with the windows.

\begin{proposition}[Connected antipodes]\label{connected-antipode}
Suppose the reconstructed bialgebra is connected, meaning $C_0=\kk1$
in the selected base category.  It has a unique charge-preserving
antipode.  Its finite restrictions and continuous extension to the
completion are compatible antipodes.  In complexes they are chain maps.
\end{proposition}

\begin{proof}
For endomorphisms define convolution by
$f*g=m(f\otimes g)\Delta$, with unit $e=\eta\varepsilon$.
Put $q=1_C-e$, so $q$ vanishes in charge zero.
In an iterated coproduct of an element of height $n$, any term with
$r>n$ factors has a height-zero factor.  Applying $q$ in every factor
kills it.  Thus
\[
 q^{*r}|_{C_\beta}=0\quad(r>\htop(\beta)),\qquad
 S=\sum_{r\geq0}(-1)^rq^{*r},\quad q^{*0}=e,
\]
is finite on each component.  Since $1_C=e+q$, cancellation of the
geometric series gives $S*1_C=1_C*S=e$.
Two-sided inverses are unique.  In height $N$ the sum stops at $r=N$,
and restriction preserves every convolution term.
The product of these compatible component maps is the continuous
completed antipode.  All factors preserve charge and parity, and in
complexes they commute with the differential.
\end{proof}

\begin{proposition}[The charge-zero Hopf condition]\label{degree-zero-hopf}
A reconstructed nonnegatively graded bialgebra is Hopf if and only if
its charge-zero bialgebra $C_0$ is Hopf.  The resulting antipodes are
compatible with finite restriction and height completion.
\end{proposition}

\begin{proof}
This is the degree-zero extension of the preceding construction
\cite[Exercise 1.4.31(b), p.~19]{GrinbergReiner}.
Let $\pi:C\to C_0$ and $\iota:C_0\to C$ be projection and inclusion.
Nonnegative grading makes both bialgebra maps.
Write $e=\eta\varepsilon$ for the convolution unit.
Given an antipode $S_0$ on $C_0$, set
\[
 a=\iota\pi,\qquad j=\iota S_0\pi,\qquad
 t=j*1_C-e.
\]
Then $j*a=a*j=e$, $t|_{C_0}=0$, and $1_C=a*(e+t)$.
The height argument in Proposition~\ref{connected-antipode} makes
$u=\sum_{r\geq0}(-1)^rt^{*r}$ locally finite and gives
$u*(e+t)=(e+t)*u=e$.  Hence $S=u*j$ satisfies
\[
 S*1_C=u*j*a*(e+t)=e,\qquad
 1_C*S=a*(e+t)*u*j=e.
\]
The construction preserves charge and the base-category structure, and
therefore restricts and extends continuously as before.
Conversely, if $S$ is an antipode on $C$, composing its two identities
with $\pi$ and $\iota$ gives the antipode $\pi S\iota$ on $C_0$.
No assumption that $S$ initially preserves charge is needed for this
necessity.  The constructed inverse and uniqueness then imply charge
preservation.
\end{proof}

Finite dimension does not imply this condition.  Concentrate
$\kk[t]/(t^2-t)$ in charge zero and set
\[
 \Delta t=t\otimes t,\qquad\varepsilon(t)=1,\qquad\Delta1=1\otimes1.
\]
These algebra maps respect $t^2=t$, and the coalgebra identities hold
on $1,t$, so this is a two-dimensional bialgebra.
An antipode would satisfy $tS(t)=1$.
Multiplication by the nonzero element $1-t$ gives a contradiction.
